% Presentation flag: uncomment to enable presentation features (slower compilation)
\newcommand*{\PRESENTATION}

% Preliminaries
\ifdefined\PRESENTATION
	\documentclass[t]{beamer}
\else
	\documentclass[t,handout]{beamer}
\fi
\mode<presentation>
\usetheme{Warsaw}
\useoutertheme{infolines}
\makeatother
\setbeamertemplate{footline}
{
  \leavevmode%
  \hbox{%
  \begin{beamercolorbox}[wd=.3\paperwidth,ht=2.25ex,dp=1ex,center]{author in head/foot}%
    \usebeamerfont{author in head/foot}\insertshortauthor
  \end{beamercolorbox}%
  \begin{beamercolorbox}[wd=.6\paperwidth,ht=2.25ex,dp=1ex,center]{title in head/foot}%
    \usebeamerfont{title in head/foot}\insertshorttitle
  \end{beamercolorbox}%
  \begin{beamercolorbox}[wd=.1\paperwidth,ht=2.25ex,dp=1ex,center]{date in head/foot}%
    \insertframenumber{} / \inserttotalframenumber\hspace*{1ex}
  \end{beamercolorbox}}%
  \vskip0pt%
}
\makeatletter
\setbeamertemplate{navigation symbols}{} 

% Includes
%\usepackage{times} % Use the times font
\usepackage{palatino} % Uncomment to use the Palatino font

\usepackage{graphicx} % Required for including images
\graphicspath{{figures/}} % Location of the graphics files
\usepackage{booktabs} % Top and bottom rules for table
\usepackage[font=small,labelfont=bf]{caption} % Required for specifying captions to tables and figures
\usepackage{wrapfig} % Allows wrapping text around tables and figures
\usepackage{lipsum,adjustbox}
\usepackage[absolute,overlay]{textpos}
\usepackage{url}
\usepackage{lmodern}
\usepackage{amsmath}
\usepackage{amsfonts}
\usepackage{color}
\usepackage{multirow}
\usepackage{multicol}
\usepackage{tikz}
\usepackage{tikz-dependency}
\usetikzlibrary{arrows.meta,graphs,graphs.standard,graphdrawing,quotes,shapes}
\usegdlibrary{layered}
\captionsetup{labelformat=empty}
\newcommand{\parser}[1]{TUPA\textsubscript{#1}}

\makeatletter
\pgfdeclareshape{vector}{
	  \inheritsavedanchors[from={rectangle}]
	  \inheritbackgroundpath[from={rectangle}]
	  \inheritanchorborder[from={rectangle}]
	  \foreach \x in {center,north east,north west,north,south,south east,south west,east,west}{
	    \inheritanchor[from={rectangle}]{\x}
	  }

    \backgroundpath{
      \pgftransformshift{\pgfpoint{-16pt}{-4pt}}
		  \draw[rounded corners=2pt] (0,0) rectangle (32pt,8pt);
    }

    \beforebackgroundpath{
      \draw[step=8pt,help lines,-] (8pt,.1pt) grid (24pt,7.9pt);
    }
}
\makeatother

% Outline slides
\AtBeginSection
{\begin{frame} \frametitle{Outline} \tableofcontents[currentsection] \end{frame}}


\begin{document}


\title{Broad-Coverage Transition-Based UCCA Parsing}
\author[Daniel Hershcovich]{\textbf{Daniel Hershcovich}$^{1,2}$ \And Omri Abend$^2$ \And Ari Rappoport$^2$
}
\institute[]{$^1$Edmond and Lily Safra Center for Brain Sciences \\
$^2$School of Computer Science and Engineering \\
Hebrew University of Jerusalem}
\date{Learning Club \\ November 17, 2016}

\begin{frame}
\titlepage

%\begin{minipage}[b]{.13\linewidth}
%\includegraphics[width=\linewidth]{huji_logo.jpg}
%\end{minipage}
%\begin{minipage}[b]{.3\linewidth}
%\includegraphics[width=.3\pagewidth]{huji_banner.png}\\
%\includegraphics[width=.4\pagewidth]{cse_banner.jpg}
%\end{minipage}
%\hfill
%\begin{minipage}[b]{.2\linewidth}
%\includegraphics[width=\linewidth]{elsc_logo.png}
%\end{minipage}
\end{frame}


%----------------------------------------------------------------------------------------

\section{Semantic Annotation Schemes}

\begin{frame}
\frametitle{Semantic Annotation Schemes}
Represent semantic structure of text as a graph.

Used by NLP applications for features and structure,
providing information such as \textit{who did what to whom?}
\vfill
Examples:
\begin{itemize}
 \item Semantic Role Labeling
 \item Semantic Dependencies
 \item Abstract Meaning Representation
 \item Universal Conceptual Cognitive Annotation
\end{itemize}
\end{frame}

\begin{frame}
\frametitle{Semantic Role Labeling (SRL)}
Annotate predicates and their arguments as a flat structure.
Examples:

\centering
\vfill
{\color{blue} PropBank}
\vfill

\begin{dependency}
	\begin{deptext}[column sep=1.5em,ampersand replacement=\^]
	After \^ graduation \^ , \^ John \^ moved \^ to \^ Paris \\
	\end{deptext}
	% PropBank
	\node[xshift=4.91em,yshift=3em,color=blue]{move.01}
	child{node at ($(\wordref{1}{5})$) {}};
	\node[xshift=-9.5em,yshift=3em,color=blue]{AM-TMP}
	child{node at ($(\wordref{1}{1})$) {}}
	child{node at ($(\wordref{1}{2})$) {}};
	\node[xshift=.65em,yshift=3em,color=blue]{A1}
	child{node at ($(\wordref{1}{4})$) {}};
	\node[xshift=10.5em,yshift=3em,color=blue]{A2}
	child{node at ($(\wordref{1}{6})$) {}}
	child{node at ($(\wordref{1}{7})$) {}};
	% FrameNet
	\node[xshift=4.91em,yshift=-3em,color=red]{Motion}
	child{node at ($(\wordref{1}{5})$) {}};
	\node[xshift=-9.5em,yshift=-3em,color=red]{Time}
	child{node at ($(\wordref{1}{1})$) {}}
	child{node at ($(\wordref{1}{2})$) {}};
	\node[xshift=.65em,yshift=-3em,color=red]{Theme}
	child{node at ($(\wordref{1}{4})$) {}};
	\node[xshift=10.5em,yshift=-3em,color=red]{Goal}
	child{node at ($(\wordref{1}{6})$) {}}
	child{node at ($(\wordref{1}{7})$) {}};
\end{dependency}

\vfill
{\color{red} FrameNet}
\end{frame}

\begin{frame}
\frametitle{Semantic Dependency Parsing (SDP)}
Graph on the text tokens, including internal structure of arguments.
Examples:

\centering
\vfill
{\color{blue} DELPH-IN MRS-derived bi-lexical dependencies (DM)}
\vfill

\begin{dependency}[theme=simple]
	\begin{deptext}[column sep=1.5em,ampersand replacement=\^]
	After \^ graduation \^ , \^ John \^ moved \^ to \^ Paris \\
	\end{deptext}
	\deproot{5}{\color{blue} top}
	\depedge{1}{2}{\color{blue} ARG2}
	\depedge{1}{5}{\color{blue} ARG1}
	\depedge{5}{4}{\color{blue} ARG1}
	\depedge{6}{5}{\color{blue} ARG1}
	\depedge{6}{7}{\color{blue} ARG2}
	\deproot[edge below]{5}{\color{red} top}
	\depedge[edge below]{5}{2}{\color{red} TWHEN}
	\depedge[edge below]{5}{4}{\color{red} ACT-arg}
	\depedge[edge below]{5}{7}{\color{red} DIR3-arg}
\end{dependency}

\vfill
{\color{red} Prague Dependency Treebank tectogrammatical layer (PSD)}
\end{frame}

\begin{frame}
\frametitle{Abstract Meaning Representation (AMR)}
Graph on knowledge resource entries inferred from the tokens.

\begin{tikzpicture}
\graph[layered layout, sibling distance=5cm, layer distance=2cm, edges={nodes={sloped}}]{
a4 Paris[as={Paris}, black];
a2 John[as={John}, black];
a1[as={person}, black];
a0[as={move-01}, black];
a3[as={city}, black];
a2[as={name}, black];
a5[as={after}, black];
a4[as={name}, black];
a6[as={graduate-01}, black];

a1 ->  ["name"' above, black] a2;
a0 ->  ["ARG0"' above, black] a1;
a0 ->  ["ARG2"' above, black] a3;
a0 ->  ["time"' above, black] a5;
a3 ->  ["name"' above, black] a4;
a2 ->  ["op1"' above, black] a2 John;
a5 ->  ["op1"' above, black] a6;
a4 ->  ["op1"' above, black] a4 Paris;
};
\draw[->, above, black] (a6) to node[sloped] {ARG0} (a1);
\end{tikzpicture}
\end{frame}

\begin{frame}
\frametitle{Universal Conceptual Cognitive Annotation (UCCA)}
\begin{tikzpicture}[level distance=25mm, sibling distance=25mm, ->,
    every circle node/.append style={fill=black}]
  \node (ROOT) [circle] {}
    child {node (After) {After} edge from parent node[left] {L}}
    child {node (graduation) [circle] {}
    {
      child {node {graduation} edge from parent node[left] {P}}
    } edge from parent node[left] {H} }
    child {node {,} edge from parent node[right] {U}}
    child {node (moved) [circle] {}
    {
      child {node (John) {John} edge from parent node[left] {A}}
      child {node {moved} edge from parent node[left] {P}}
      child {node [circle] {}
      {
        child {node {to} edge from parent node[left] {R}}
        child {node {Paris} edge from parent node[right] {C}}
      } edge from parent node[right] {A} }
    } edge from parent node[right] {H} }
    ;
  \draw[dashed,->] (graduation) to node [auto] {A} (John);
\end{tikzpicture}
\begin{wraptable}{l}{4.8mm}
  \vspace{-3cm}
  \scalebox{.75}{
  \begin{adjustbox}{margin=2mm,frame}
  \begin{tabular}{cl|cl|cl}
	  P & process &
	  S & state &
	  A & participant \\
	  L & linker &
	  H & linked scene &
	  C & center \\
	  E & elaborator &
	  D & adverbial &
	  R & relator \\
	  N & connector &
	  U & punctuation &
	  F & function unit \\
	  G & ground &&
	  \multicolumn{3}{c}{- - - remote edge}
  \end{tabular}
  \end{adjustbox}
  }
\end{wraptable}
\end{frame}

\begin{frame}
\frametitle{Universal Conceptual Cognitive Annotation (UCCA)}
\begin{wraptable}{r}{35mm}
  \vspace{-5mm}
  \begin{adjustbox}{margin=2mm,frame}
  \begin{tabular}{cl}
	  P & process \\
	  S & state \\
	  A & participant \\
	  L & linker \\
	  H & linked scene \\
	  C & center \\
	  E & elaborator \\
	  D & adverbial \\
	  R & relator \\
	  N & connector \\
	  U & punctuation \\
	  F & function unit \\
	  G & ground
  \end{tabular}
  \end{adjustbox}
  \captionof{table}{Edge labels of the foundational layer.}
\end{wraptable}
Semantic representation scheme giving an alternative approach to syntax.
\begin{itemize}
 \item Builds on typological and Cognitive Linguistics literature \cite{Dixon:basic,croft2004cognitive}.
 \item Demonstrated applicability to English, French, German \& Czech.
 \item Support for rapid annotation.
 \pause
 \item Semantic stability in translation \cite{sulem2015conceptual}.
 \item Useful for machine translation evaluation \cite{birch2016hume}.
 \item Applicability limited by absence of parser.
 \pause
 \item Also supported: linkage (discourse relations) and implicit units,
 but deferred to later work.
\end{itemize}
\end{frame}

\begin{frame}
\frametitle{Structural Properties}
\noindent
(1) non-terminal nodes, (2) reentrancy, (3) discontinuity
\centering
\vfill
\begin{minipage}{.48\linewidth}{\centering
  \begin{tikzpicture}[level distance=12mm, sibling distance=16mm, ->,
      every node/.append style={midway}]
    \node (ROOT) [fill=black, circle] {}
      child {node [fill=black, circle] {}
      {
        child {node {John} edge from parent node[left] {\scriptsize $C$}}
        child {node {and} edge from parent node[left] {\scriptsize $N$}}
        child {node {Mary} edge from parent node[left] {\scriptsize $C$}}
      } edge from parent node[left] {\scriptsize $A$} }
      child {node {went} edge from parent node[left] {\scriptsize $P$}}
      child {node {home} edge from parent node[left] {\scriptsize $A$}}
      ;
  \end{tikzpicture}
  }
\end{minipage}
\hfill
\begin{minipage}{.48\linewidth}{\centering
  \begin{tikzpicture}[level distance=12mm, sibling distance=2cm, ->,
      every node/.append style={midway}]
    \node (ROOT) [fill=black, circle] {}
      child {node {John} edge from parent node[left] {\scriptsize $A$}}
      child {node [fill=black, circle] {}
      {
      	child {node {gave} edge from parent node[left] {\scriptsize $C$}}
      	child {node (everything) {everything} edge from parent[white]}
      	child {node {up} edge from parent node[left] {\scriptsize $C$}}
      } edge from parent node[right] {\scriptsize $P$} }
      ;
    \draw[bend right,->] (ROOT) to[out=-20, in=180] node [left] {\scriptsize $A$} (everything);
  \end{tikzpicture}
  }
\end{minipage}

\vspace{-1mm}
  \begin{tikzpicture}[level distance=12mm, sibling distance=2cm, ->,
      every node/.append style={midway}]
    \node (ROOT) [fill=black, circle] {}
      child {node (John) {John} edge from parent node[left] {\scriptsize $A$}}
      child {node {decided} edge from parent node[left] {\scriptsize $P$}}
      child {node (totakeaquickshower) [fill=black, circle] {}
      {
        child {node {to} edge from parent node[left] {\scriptsize $F$}}
        child {node (takeashower) [fill=black, circle] {}
        {
          child {node {take} edge from parent node[left] {\scriptsize $C$}}
          child {node {a} edge from parent node[right] {\scriptsize $F$}}
          child {node (quick) {quick} edge from parent[white]}
          child {node {shower} edge from parent node[right] {\scriptsize $C$}}
        } edge from parent node[right] {\scriptsize $P$} }
      } edge from parent node[left] {\scriptsize $A$} }
      ;
    \draw[bend left,dashed,->] (takeashower) to node [auto] {\scriptsize $A$} (John);
    \draw[bend left,->] (totakeaquickshower) to node [auto] {\scriptsize $D$} (quick);
  \end{tikzpicture}
\end{frame}


\begin{frame}
\frametitle{UCCA Corpora}
\centering
\begin{tabular}{l|ccc|c}
& \multicolumn{3}{c|}{Wiki} & 20K \\
& \small Train & \small Dev & \small Test & Leagues \\
\hline
\# passages & 300 & 34 & 34 & 154 \\
\# sentences & 4267 & 453 & 518 & 506 \\
\hline
\# nodes & 298,665 & 33,263 & 37,262 & 29,315 \\
\% terminal & 42.95 & 43.62 & 42.89 & 42.09 \\
\% non-term. & 58.30 & 57.46 & 58.31 & 60.01 \\
\% \bf discont. & \bf 0.53 & \bf 0.51 & \bf 0.47 & \bf 0.81 \\
\% \bf reentrant & \bf 2.31 & \bf 1.76 & \bf 2.18 & \bf 2.03 \\
\hline
\# edges & 287,381 & 32,015 & 35,846 & 27,749 \\
\% primary & 98.29 & 98.81 & 98.75 & 97.73 \\
\% remote & 1.71 & 1.19 & 1.25 & 2.27 \\
\hline
\multicolumn{3}{l}{\footnotesize Average per non-terminal node} \\
\# children & 1.67 & 1.68 & 1.66 & 1.61 
\end{tabular}
\captionof{table}{Corpus statistics
excluding root, implicit and linkage units.}
\end{frame}



\section{Transition-based UCCA Parsing}

\begin{frame}
\frametitle{Transition-Based Parsing}
\begin{itemize}
 \item Parse sentence $w_1 \ldots w_n$ to graph $G=(V,E,\ell)$ incrementally.
 \item Using buffer $B$ and stack $S$.
 \item Classifier determines transition to apply at each step.
 \item Trained by an oracle based on gold-standard graph.
\end{itemize}

\begin{center}
	\begin{tikzpicture}
	\draw[xstep=1cm,ystep=5mm,color=gray] (-0.01,0) grid (1,.5);
	\node[anchor=west] at (.2,1) {$S$};
	\node[fill=black, circle] at (.5,.25) {};
	\draw[xstep=18mm,ystep=5mm,color=gray] (1.79,0) grid (12.4,.5);
	\node[anchor=west] at (1.9,1) {$B$};
	\node[anchor=west] at (2,.25) {After};
	\node[anchor=west] at (3.5,.25) {graduation};
	\node[anchor=west] at (5.6,.25) {John};
	\node[anchor=west] at (7.2,.25) {moved};
	\node[anchor=west] at (9.4,.25) {to};
	\node[anchor=west] at (10.9,.25) {Paris};
	\end{tikzpicture}
\end{center}

\vfill
\pause
Transitions defined for UCCA parsing:\\
\{\textsc{Shift, Reduce, Node$_X$, Left-Edge$_X$, Right-Edge$_X$,}\\
\hspace{5mm}\textsc{Left-Remote$_X$, Right-Remote$_X$, Swap, Finish}\}

\vfill
\pause
Supports non-terminal nodes, reentrancy and discontinuity.
\end{frame}

\begin{frame}
\frametitle{Transition-Based Parsing}
\begin{center}
	\begin{tikzpicture}[level distance=15mm, sibling distance=2cm]
	\draw[xstep=1cm,ystep=5mm,color=gray] (-0.01,0) grid (3,.5);
	\node[anchor=west] at (1.2,1) {$S$};
	\node[fill=black, circle] at (.5,.25) {};
	\node[fill=blue, circle] at (1.5,.25) {};
	\node[anchor=west] at (2,.25) {John};
	\draw[xstep=12mm,ystep=5mm,color=gray] (3.59,0) grid (7.2,.5);
	\node[anchor=west] at (5,1) {$B$};
	\node[anchor=west] at (3.5,.25) {moved};
	\node[anchor=west] at (5,.25) {to};
	\node[anchor=west] at (6,.25) {Paris};
	\node[anchor=west] at (8,1) {$G$};
	\node[fill=black, circle] at (9,.5) {}
	  child {node  {\small After} edge from parent [->] node[left] {\small L}}
	  child {node [fill=blue, circle] {}
	  {
	    child {node {\small graduation} edge from parent [->] node[right] {\small P}}
	  } edge from parent [->] node[right] {\small H} };
	\node[anchor=west] at (0,-1) {History:};
	\node[anchor=west] at (1.5,-1) {\textsc{Shift},};
	\node[anchor=west] at (1.5,-1.5) {\textsc{Right-Edge\textsubscript L},};
	\node[anchor=west] at (1.5,-2) {\textsc{Reduce},};
	\node[anchor=west] at (1.5,-2.5) {\textsc{Shift},};
	\node[anchor=west] at (1.5,-3) {\textsc{Node\textsubscript P},}; 
	\node[anchor=west] at (1.5,-3.5) {\textsc{Reduce},};
	\node[anchor=west] at (1.5,-4) {\textsc{Shift},};
	\node[anchor=west] at (1.5,-4.5) {\textsc{Right-Edge\textsubscript H},};
	\node[anchor=west] at (1.5,-5) {\textsc{Shift}};
	\end{tikzpicture}
	\captionof{figure}{Example for intermediate state during transition-based parsing.}
\end{center}
\end{frame}

\begin{frame}
\frametitle{\parser{} Classifiers}
We experiment with three classifiers:
\begin{flushleft}
	\begin{tabular}{ll}
	\textbf{\parser{sparse}} & Perceptron, sparse features: words, POS tags \& edge label combinations. \\
	\textbf{\parser{dense}} & Perceptron, dense embedding features: word2vec \cite{mikolov2013efficient} for words, else random. \\
	\textbf{\parser{NN}} & 2-layer MLP, learned embedding features, logistic activation + dropout. \\
	\end{tabular}
\end{flushleft}
\vspace{5mm}
For all classifiers, inference is performed greedily, i.e., without beam search.
\end{frame}

%\begin{frame}
%\frametitle{\parser{} Features}
%\begin{adjustbox}{margin=3pt,frame}
%\begin{tabular}{l}
%{\footnotesize Features from \cite{zhang2009transition}:} \\
%\textbf{unigrams} \\
%$s_0te, s_0we, s_1te, s_1we, s_2te, s_2we, s_3te, s_3we,$ \\
%$b_0wt, b_1wt, b_2wt, b_3wt,$ \\
%$s_0lwe, s_0rwe, s_0uwe, s_1lwe, s_1rwe, s_1uwe$ \\
%\textbf{bigrams} \\
%$s_0ws_1w, s_0ws_1e, s_0es_1w, s_0es_1e, s_0wb_0w, s_0wb_0t,$ \\
%$s_0eb_0w, s_0eb_0t, s_1wb_0w, s_1wb_0t, s_1eb_0w, s_1eb_0t,$ \\
%$b_0wb_1w, b_0wb_1t, b_0tb_1w, b_0tb_1t$ \\
%\textbf{trigrams} \\
%$s_0es_1es_2w, s_0es_1es_2e, s_0es_1eb_0w, s_0es_1eb_0t,$ \\
%$s_0es_1wb_0w, s_0es_1wb_0t, s_0ws_1es_2e, s_0ws_1eb_0t$ \\
%\textbf{separator} \\
%$s_0wp, s_0wep, s_0wq, s_0wcq, s_0es_1ep, s_0es_1eq,$ \\
%$s_1wp, s_1wep, s_1wq, s_1weq$ \\
%
%\textbf{extended} \footnotesize \cite{zhu2013fast} \\
%$s_0llwe, s_0lrwe, s_0luwe, s_0rlwe, s_0rrwe,$ \\
%$s_0ruwe, s_0ulwe, s_0urwe, s_0uuwe, s_1llwe,$ \\
%$s_1lrwe, s_1luwe, s_1rlwe, s_1rrwe, s_1ruwe$ \\
%\end{tabular}
%\begin{tabular}{l}
%\textbf{disco} \footnotesize \cite{maier2015discontinuous} \\
%$s_0xwe, s_1xwe, s_2xwe, s_3xwe,$ \\
%$s_0xte, s_1xte, s_2xte, s_3xte,$ \\
%$s_0xy, s_1xy, s_2xy, s_3xy$ \\
%$s_0xs_1e, s_0xs_1w, s_0xs_1x, s_0ws_1x, s_0es_1x,$ \\
%$s_0xs_2e, s_0xs_2w, s_0xs_2x, s_0ws_2x, s_0es_2x,$ \\
%$s_0ys_1y, s_0ys_2y, s_0xb_0t, s_0xb_0w$ \\
%
%{\footnotesize Features from \cite{tokgoz2015transition}:} \\
%\textbf{counts} \\
%$s_0P, s_0C, s_0wP, s_0wC, b_0P, b_0C, b_0wP, b_0wC$ \\
%\textbf{edges} \\
%$s_0s_1, s_1s_0, s_0b_0, b_0s_0, s_0b_0e, b_0s_0e$ \\
%\textbf{history} \\
%$a_0, a_1$ \\
%
%\textbf{remote} \footnotesize (Novel, UCCA-specific features) \\
%$s_0R, s_0wR, b_0R, b_0wR$
%\end{tabular}
%\end{adjustbox}
%\captionsetup{singlelinecheck=off}
%\captionof{table}{
%  Binary feature templates for \parser{sparse}. Notation:\\
%  $s_i$, $b_i$: $i$th stack and buffer items.\\
%  $w$, $t$: word form and POS tag of the terminal returned by $h^*(\cdot)$
%  (see Appendix~\ref{appendix:conversion}).\\
%  $e$: edge label to the node returned by $h(\cdot)$.\\
%  $l$, $r$ ($ll$, $rr$): leftmost and rightmost (grand)children.\\
%  $u$ ($uu$): unary (grand)child, when only one exists.\\
%  $p$: unique separator punctuation between $s_0$ and $s_1$. $q$: separator count.\\
%  $x$: gap type (``none'', ``pass'' or ``gap'') at the sub-graph under the current node.\\
%  $y$: sum of gap lengths \protect\cite{Maier:Lichte:11}.\\
%  $P$, $C$: number of parents and children.\\
%  $R$: number of remote children.\\
%  $a_i$: action taken $i$ steps back.
%}
%
%We also employ a real-valued feature,
%\textbf{ratio}, corresponding to the ratio between the number of terminals to number of nodes
%in the graph $G$.
%This novel feature serves as a regularizer for the creation of new nodes,
%and should be beneficial for other transition-based constituency parsers too.
%\end{frame}








\section{Experiments}

\begin{frame}
\frametitle{Experimental Setup}
We conduct our main experiment on the UCCA Wikipedia corpus,
and use the English part of the UCCA \textit{Twenty Thousand Leagues Under the Sea}
English-French parallel corpus as out-of-domain data.
\end{frame}

\begin{frame}
\frametitle{Evaluation}
We report two variants of labeled precision, recall and F-score:
one where we consider only primary edges, and another for remote edges.
Given graphs $G_p=(V_p,E_p,\ell_p)$ and $G_g=(V_g,E_g,\ell_g)$
over terminals $W = \{w_1,\ldots,w_n\}$,
the yield $y(e) \subseteq W$ of an edge $e=(u,v)$ in either graph
is the set of terminals in $W$ that are descendants of $v$.
The \textit{mutual edges} between the graphs are:
\begin{flushleft}
\begin{tabular}{lcr}
	\multicolumn{3}{l}{$M(G_p,G_g) =
    \left\{(e_1,e_2) \in E_p \times E_g \;|\;
    y(e_1) = y(e_2) \wedge \ell_p(e_1)=\ell_g(e_2)\right\}$} \vspace{5mm} \\
    and we define\\
	$\text{LP} = |M(G_p,G_g)| / |E_p|$ \hspace{2cm} &
	$\text{LR} = |M(G_p,G_g)| / |E_g|$ \hspace{2cm} &
	$\text{LF} = 2 \cdot \text{LP} \cdot \text{LR} / (\text{LP} + \text{LR})$.
\end{tabular}
\end{flushleft}
\end{frame}

\begin{frame}
\frametitle{Baselines}
\begin{wrapfigure}[11]{r}{14cm}
	\centering
	\scalebox{.7}{
	\begin{dependency}[theme = simple]
	\begin{deptext}[column sep=.7em,ampersand replacement=\^]
	After \^ graduation \^ , \^ John \^ moved \^ to \^ Paris \\
	\end{deptext}
	\depedge{2}{1}{L}
	\depedge{2}{3}{U}
	\depedge[dashed]{2}{4}{A}
	\depedge{5}{4}{A}
	\depedge{2}{5}{H}
	\depedge{7}{6}{R}
	\depedge{5}{7}{A}
	\end{dependency}
	}
	\scalebox{.7}{
	\begin{dependency}[theme = simple]
	\begin{deptext}[column sep=.7em,ampersand replacement=\^]
	John \^ gave \^ everything \^ up \\
	\end{deptext}
	\depedge{1}{2}{A}
	\depedge{3}{2}{A}
	\depedge{4}{2}{C}
	\end{dependency}
	}
	\scalebox{.7}{
	\begin{dependency}[theme = simple]
	\begin{deptext}[column sep=.7em,ampersand replacement=\^]
	John \^ and \^ Mary \^ went \^ home \\
	\end{deptext}
	\depedge[edge start x offset=-6pt]{1}{4}{A}
	\depedge{2}{1}{N}
	\depedge{3}{1}{C}
	\depedge{5}{4}{A}
	\end{dependency}
	}
	\captionof{figure}{Bilexical approximation for UCCA graphs.}
\end{wrapfigure}

Since there are no existing UCCA parsers, we use bilexical DAG parsers:
\begin{enumerate}
 \item Convert UCCA into bilexical dependencies.
 \item Train parsers on the resulting training set.
 \item Apply trained parsers to the test set.
 \item Reconstruct UCCA graphs.
 \item Compare with gold standard.
\end{enumerate}
Upper bounds are computed by applying
the conversion both ways to the gold standard
graphs and comparing to the original.
\end{frame}

\begin{frame}
\frametitle{Results}
\parser{NN} obtains the highest scores in nearly all metrics:
	  
\begin{center}
	\begin{tabular}{l|ccc|ccc||ccc|ccc}
	& \multicolumn{6}{c||}{Wiki (in-domain)} & \multicolumn{6}{c}{20K Leagues (out-of-domain)} \\
	& \multicolumn{3}{c|}{Primary} & \multicolumn{3}{c||}{Remote}
	& \multicolumn{3}{c|}{Primary} & \multicolumn{3}{c}{Remote} \\
	& \textbf{LP} & \textbf{LR} & \textbf{LF} & \textbf{LP} & \textbf{LR} & \textbf{LF}
	& \textbf{LP} & \textbf{LR} & \textbf{LF} & \textbf{LP} & \textbf{LR} & \textbf{LF} \\
	\hline
	\multicolumn{4}{l}{\rule{0pt}{2ex} \footnotesize Bilexical Approximation} \\
	Upper Bound % on the test set and ood set only
	& 93.4 & 83.7 & 88.3 & 73.9 & 49.5 & 59.3
	& 93.5 & 83.5 & 88.2 & 66.7 & 31.6 & 42.9 \vspace{.1cm} \\
	DAGParser \cite{ribeyre-villemontedelaclergerie-seddah:2014:SemEval}
	& 63.7 & 	56.1	 & 59.5	 & 0.8	 & 9.5	 &  1.4
	& 58	 & 49.8	 & 53.4 & -- & 0 & 0 \\
	TurboParser \cite{almeida-martins:2015:SemEval}
	& 60.2	 & 47.4	 & 52.9	 & 2.2	 & 7.8	 &  3.4
	& 52.6	 & 39	 & 44.7	 & 100	 & 0.3	 & 0.6 \\
	\hline
	\multicolumn{4}{l}{\rule{0pt}{2ex} \footnotesize Direct Approach} \\
	\parser{sparse}
	& 64 & 55.6 & 59.5 & 16 & 11.6 & 13.4 
	& 60.6 & 53.9 & 57.1 & 20.2 & 10.3 & 13.6 \\
	\parser{dense} 
	& 55 & 54.8 & 54.9 & 15.2 & {\bf 16.9} & {\bf 16}
	& 54.8 & 55.2 & 55 & 6 & 3 & 4 \\
	\parser{NN}
	& {\bf 65} & {\bf 62.5} & {\bf 63.7} & {\bf 20.7} & 11.3 & 14.6
	& {\bf 58.3} & {\bf 56.4} & {\bf 57.3} & {\bf 15.2} & {\bf 3.8} & {\bf 6}
	\end{tabular}
\end{center}
\end{frame}

\begin{frame}
\frametitle{Tree Approximation}
For completeness, we also explore lossily converting UCCA into trees,
resulting in a simplified task for the underlying parser,
in addition to the maturity of tree-based parsers.

Although remote edges are of pivotal importance, exploring tree approximation methods
can inform the future development of DAG parsers in general and of UCCA parsers in particular.
\begin{wraptable}[11]{l}{15cm}
\hspace{-6mm}
\begin{tabular}{l|ccc}
\hline
\multicolumn{4}{l}{\rule{0pt}{2ex} \footnotesize Constituency Tree Approximation} \\
Upper Bound & 100 & 100 & 100 \vspace{.1cm} \\
\textsc{uparse} \cite{maier-lichte:2016:DiscoNLP} & 63 & 64.7 & 63.7 \\
\hline
\multicolumn{4}{l}{\rule{0pt}{2ex} \footnotesize Dependency Tree Approximation} \\
Upper Bound & 93.7 & 83.6 & 88.4 \vspace{.1cm} \\
MaltParser \cite{nivre2007maltparser} & 64.9 & 57.9 & 61 \\
LSTM Parser \cite{dyer2015transition} & {\bf 74.9} & {\bf 66.4} & {\bf 70.2} \\
\hline
\multicolumn{4}{l}{\rule{0pt}{2ex} \footnotesize Direct Tree Parsing} \\
\parser{sparse} $-$ \textsc{Remote} & 65.5 & 57.5 & 61.3 \\
\parser{dense} $-$ \textsc{Remote} & 57.2 & 57.3 & 57.2 \\
\parser{NN} $-$ \textsc{Remote} & 66.3 & 64.4 & 65.3 \\
\end{tabular}
\end{wraptable}

The performance is encouraging in light of
UCCA's inter-annotator agreement of 80--85\%
F-score on primary edges \cite{abend2013universal}.
\end{frame}



\section{Future Work}

\begin{frame}
\frametitle{UCCA-Based Distributed Representation}
Vector representation for sentences and documents,
based on recursive composition on the UCCA graph.

\vfill
Impact:
\begin{itemize}
\item General automatic semantic feature extractor for text.
\item Accurate measure for text similarity.
\item Understand the semantic contribution of different elements.
\end{itemize}

\vfill
\begin{center}
  \begin{tikzpicture}[level distance=1cm, sibling distance=2cm, ->]
    \node (ROOT) [fill=black, vector] {}
      child {node (After) {After} edge from parent [shorten <=3mm]
        node[left] {\scriptsize $L$\;}}
      child {node (graduation) [fill=black, vector] {}
      {
        child {node {graduation} edge from parent node[left] {\scriptsize $P$}}
      } edge from parent node[left] {\scriptsize $H$} }
      child {node {,} edge from parent node[right] {\scriptsize $U$}}
      child {node (moved) [fill=black, vector] {}
      {
        child {node (John) {John} edge from parent [shorten <=2mm]
          node[left] {\scriptsize $A$}}
        child {node {moved} edge from parent node[left] {\scriptsize $P$}}
        child {node [fill=black, vector] {}
        {
          child {node {to} edge from parent node[left] {\scriptsize $R$}}
          child {node {Paris} edge from parent node[left] {\scriptsize $C$}}
        } edge from parent [shorten <=2.5mm,shorten >=3mm] node[left] {\scriptsize $A$} }
      } edge from parent [shorten <=3.5mm,shorten >=4mm] node[right] {\scriptsize $H$} }
      ;
    \draw[dashed,->,shorten <=2.5mm] (graduation) to node [auto] {\scriptsize $A$} (John);
    \node (LKG) at (-1.8,0) [fill=black!20, vector] {};
        \draw[bend right,shorten <=4mm] (LKG) to
          node [auto, left] {\scriptsize $LR$} (After);
        \draw (LKG) to[out=-60, in=120]
          node [below] {\scriptsize $LA$\quad\;} (graduation);
        \draw[shorten <=2mm] (LKG) to[out=30, in=90]
          node [above] {\scriptsize $LA$} (moved);
  \end{tikzpicture}
\end{center}
\end{frame}




\section{Conclusions}

\begin{frame}
\frametitle{Conclusions}
We present \parser{}, the first parser for UCCA, and
evaluate it in both in-domain and out-of-domain settings,
showing it surpasses bilexical DAG parsers on the task of UCCA parsing.

Future work will incorporate LSTMs into \parser{},
and apply the parser to more languages such as German,
demonstrating the importance of broad-coverage parsing.
We will also improve the conversion-based methods and explore different target representations.
A UCCA parser will enable using the scheme for representation in NLP tasks.

\begin{itemize}
 \item UCCA exhibits formal properties important for semantic representation.
 \item We present the first parser for UCCA and the first to support these properties.
\end{itemize}

Corpora are available at \url{http://www.cs.huji.ac.il/~oabend/ucca.html}.
Parser code available at \url{https://github.com/danielhers/ucca}.
Written in Python using Keras: \url{https://github.com/fchollet/keras}.
\end{frame}

\begin{frame}
\vfill
\begin{center}
Thank you
\end{center}
\end{frame}



\begin{frame}[allowframebreaks]
\frametitle{References}
\bibliographystyle{apalike}
\tiny\bibliography{references}
\end{frame}

\end{document}
